\section{Heard Around Town}

\begin{quotex}
Arseny's soul wanted to touch Ursina's soul ... Get used to separation, said Death, it is painful, even if it is only temporary. Will we recognize each other in eternity? asked Arseny's soul. That depends in large part on you, said Death: souls often harden during the course of life and then they barely recognize anyone after death.''
\flright{\textsc{Eugene Vodolazkin}, \emph{from Laurus}}
\end{quotex}


\paragraph{Potter's Field}


\begin{quotex}
The potter's field was a mournful place... There lay the plague dead, pilgrims, the strangled, unbaptized babies, and suicides; those drowned by waters, taken by battle, murdered, and stricken by fire. Suddenly surprised, those who had fallen from lightning, dead from frost and every sort of wound. The lives of those unfortunates were varied and it was not life that united them: their resemblance to one another consisted of death. It was death without Confession.
\flright{\textsc{Eugene Vodolazkin}, \emph{from Laurus}}
\end{quotex}


That was the situation in Russia, some six centuries ago. Potter's field was an open pit into which the corpses were dropped. Once a year, a priest gave a blessing, the pit was filled in, and another one dug.

If sudden death without benefit of the last sacraments is a curse, then presumably a slow death, preceded by a gradual breakdown of organs, must be a blessing.

A few years ago I was sitting across the desk of a specialist. He looked at my medical history and test results, then looked up at me. He read another page and looked up again. Finally, after he had read the full report, he said to me, ``Well, at least you \emph{look} good!''

\paragraph{Mohammed and Charlemagne}


I made \textbf{Julius Evola}'s review of \emph{Mohammed and Charlemagne}\footnote{\url{https://gornahoor.net/?p=15799}} available for its possible relevance today. \textbf{Henri Pirenne}'s thesis has been influential, even today I believe. As Evola points out, it has the defect of explaining everything in terms of material causes: economic activity, trade routes, etc. There are several things to notice in the review.

\textbf{Pace materialists}: Attention must be paid to the difference between necessary and sufficient causes\footnote{\url{https://en.wikipedia.org/wiki/Necessity_and_sufficiency}}. Even if certain material circumstances apply pressure to social structures, that alone cannot explain everything. Specifically, material forces are the passive elements, leaving out the active elements. Agency must be accounted for and cannot be explained in material terms.

\textbf{Pace neo-pagans}: The Holy Roman Empire was Germanic, true Indo-European civilization. Evola used the word ``Aryan'', but that word has changed its meaning since 1939. Rather, we can understand Indo-European in the way used by \textbf{Georges Dumezil\footnote{\url{https://gornahoor.net/?p=5198}}}. Today, we often hear of saving ``white'' or ``European'' civilization, which is an imprecise way of expressing it. This is what is meant by that:

\begin{quotex}
The Middle Ages took on its own appearance and which, in its great political-social structures, reproduced essentially the distinctive traits of all the Indo-European civilizations of Antiquity
\end{quotex}


So to be precise: there must first be a tradition to save. ``Europe'' as a geographical designation or ``white'' as a biological category, are simply forms of materialism and would suffer from the same explanatory defects mentioned by Evola about Pirenne.

That Tradition, as Evola asserts, was the Catholic religion (whatever you may think of its current manifestation). That religion was the bearer of an older tradition. Evola explains:

\begin{quotex}
We find these traditions i.e., the Germanic in a state of involution and degeneration in the course of the period that we have called the interregnum, but not at the point of preventing, under certain conditions, a resumption, an awakening, a galvanization. That is exactly what happened, thanks to the tradition already mentioned and to a certain heroic morality, when an historic conjuncture led to the formation of the Holy Roman Empire. Here, the Roman symbol was the common reference of a \emph{supernatural order proposed by Christianity, which acted, so to say, as a catalyst on the manifestation and the survival of latent spiritual and racial forces}.
\end{quotex}


There was no mistake, no deception as neo-pagans claimed. To the contrary, they don't even seem to know what that latent tradition was, in terms of its understanding of the supernatural order, social order, and even a heroic morality.

\textbf{Pace capitalists}: Capitalists and leftists regard Medieval society as regressive and decadent. In the communist historical scheme, capitalism replaced feudalism and communism will replace capitalism. Hence, capitalism is just one stage of the revolutionary mindset.

Evola asserts, on the contrary, that the disdain for mercantilism represents a higher value. Hence, the agrarian and feudal system, centered on the transcendent, represents a higher value.

\textbf{Pace Evola} (perhaps): It is unclear how Evola regarded the influence of the transcendent, or perhaps he regarded it all as the expression of transcendent values. In other words, a genuine Tradition not only has a theory of the supernatural order, it also provides access to that order. The third dimension of history, then, implies divine interference.

Therefore, a Traditional society should be the beneficiary of divine blessings. Hence, the emergence of Europe, from the time of Charlemagne as long as it was Christian prospered. The French revolution, which inaugurated a new worldview was the overthrowing of Christianity. That process was completed with the end of last remnants of the Holy Roman Empire last century. Since then, Europe has been in turmoil, with periods of death and destruction, with what seems to be a lack of the will to live.

This is the result of the re-paganization of Europe, which has failed to catalyze ``a resumption, an awakening, a galvanization'' of traditional values. That is simply a fact. Pagans will consider this condition to be the will of the gods.

\paragraph{Information, Understanding, and Esoterism}

\begin{quotex}
I have guarded in my memory two precious treasures of knowledge which I received from the Messenger of God. One of them I have made public, but if I were to divulge the other, you would cut my throat.
\flright{\textsc{Abu Hurayrah}}
\end{quotex}


Information is the accumulation of facts, scholarly documentation, and so on. Without understanding, that is, without the ability to grasp the accumulation as a ``whole'', there is only effete erudition. It is necessary to learn to think like the angels, who understand a myriad of specific facts because of the understanding of a paucity of higher principles.

Here are some steps to reach understanding:

\begin{itemize}


\item Follow logical consequences of a thesis, even if not explicitly affirmed

\item Put ideas into a larger whole,

\item Bring out the inner meaning of ideas
\end{itemize}


\textbf{Frithjof Schuon} explains the last step, which may require some thought:

\begin{quotex}
Esoterism refers to the nature of things, hence to intrinsic, and not legal or conventional, values.
\end{quotex}


These days, one hears a lot about ``tradition'', whether such and such idea or some so and so is ``traditional''. The first question to ask is: ``Can you build a civilization on it?'' For example, can you build a society on the basis of a man-made myth of Cthulhu\footnote{\url{https://en.wikipedia.org/wiki/Cthulhu}}? Some seem to think so.

\paragraph{Putting on a Show}


Someone asked me recently about whether Guenon was ``putting on a show'' with his conversion to Islam. A rather strange question, since who would you be when the show ends? To put on a really big show\footnote{\url{https://www.youtube.com/watch?v=xnNM-vZoaHw}}, you have to go all in and not pretend.

\paragraph{Gangs of Feral Men}


Another distortion of Tradition comes from a new movement calling itself neo-reaction. Rather than being formed by years of immersion in the likes of \textbf{Edmund Burke} and \textbf{Joseph de Maistre}, and their many followers and expositors, they read a web site last month.

A recent result is the claim that bands of feral men\footnote{\url{http://www.socialmatter.net/2016/02/23/mannerbund-101/}} were the source of civilization, rather than the threat they really are. In the face of all of history, common sense, and higher thinking, he asserts that the family is not the foundation of society. There is no understanding of ``organic'', since it is asserted that it is the result of some ``male competitive instinct''. Rather, organic is intrinsic and essential; that is, is cannot be the result of something more fundamental, since it is fundamental. A group of random men getting together for a contingent purpose is just the opposite of ``organic''.

A foundation of a real civilization is not some pseudo-Nietzschean will to power, but is rather its transcendent source. Thus, according to Evola, the Goths called their kings ``divine heroes'' and they believed their origin was divine, not instinctual.

The essay assumes some strange update of Rousseau's Social Contract and asserts that women are the ``property of men''. If you decide to read that essay and want to restore your sanity, then you can start with Family values\footnote{\url{https://gornahoor.net/?p=4534}}

And for an alternative view of the origins of civilization, try Where Have All the Heroes Gone?\footnote{\url{https://gornahoor.net/?p=6067}}

Of course, the thesis that men will naturally arrange themselves in a hierarchy can be easily proved, when the intellectually inferior begin to defer to their intellectual superiors.

\paragraph{Genetics and the Gods}


I read one comment about the existence of some deviant behavior. It \emph{must} have some survival value, the comment asserted. But that is precisely the pagan view: whatever is, must have been willed by the gods.

Obviously that is both bad science and bad theology. Whatever genetic features (assuming even that it is genetic) exists is not necessarily selected for. Rather it could be mutations that have not been eradicated. Deleterious mutations in the genetic code, or ``genetic load'', are always arising. Moreover, there may be group selections: a disordered society may select for various psychological disorders.

True Providence is the Realm of Being, of possibilities of manifestation, the world of perfection. That is what God wills. The Realm of Existence, or Becoming, is less certain. Therefore, it requires discernment to determine the will of God in any situation.

\paragraph{Detachment from the World is Attachment to God}


I listened to an online sermon with that title.\footnote{\textit{Detachment from the World is Attachment to God}: \url{https://www.youtube.com/watch?v=4Oo8HO1obOo\&list=PLtdt1LQlokf-evtBjgotlYNnBXLJo-lhC\&index=2}}

Thanks to youtube matching, it pulled up a lecture by a yogi with the same title. It was not in English, so I don't know what was said. I assume the point was the same.

\paragraph{The Last Judgment}


\begin{quotex}
He sitteth on the right hand of the Father, God Almighty, from whence he shall come to judge the quick and the dead. At whose coming all men shall rise again with their bodies, and \emph{shall give account for their own works}.
\flright{\textit{from the Athanasian Creed}}

Before this throne, in my vision, the dead must come, great and little alike; and the books were opened. Another book, too, was opened, the book of life. And the dead were judged by their deeds, as the books recorded them.
\flright{\textsc{Revelation 20:12}}

The just shall appear whiter, more beautiful, and more resplendent than the sun. ... on the other hand, the bodies of the reprobate shall appear black and hideous, and shall send forth an intolerable stench.
\flright{\textsc{St. Alphonsus Liguori}, \emph{Preparation for Death}}
\end{quotex}


According to the saints, the general judgment is the third coming of Christ, not the second that is commonly supposed. At that time, all souls will be exposed: some will be pure, most will be hardened and barely recognizable.

The purpose is to demonstrate the justice of the judgment. The consequences of one's thoughts, words, and deeds will be fully revealed. In our current condition, we can seldom see them beyond a few people close to us or a few years in the future. However, in the general judgment, we will see how our thoughts, words, and deeds rippled through the world and down through the generations. One bad choice now may be felt for several generations.

Our innermost thoughts will be exposed. Perhaps you lusted after your neighbor's wife, or worse, his daughter, or even worse, your own step-daughter. Maybe you gossiped about or harshly criticized a friend or relative behind his back; that will all become known. You get the picture. The point is to prepare for it by keeping your soul pure right now.

You say you don't believe in such superstitions. Does that really make it OK? Do you really mean that you can lust, hate, steal, or bring pain just because no one will ever find out? That is really sad. You would be better off believing the superstition than believing your own denials, lies, and self-deception.


\flrightit{Posted on 2016-02-28 by Cologero}

\begin{center}* * *\end{center}

\begin{footnotesize}
\begin{sffamily}
\texttt{Matt on 2016-02-29 at 06:39 said:}

It's interesting to see people who would balk at John Calvin's thought all of a sudden became unacknowledged Calvinists with respect to the divine will. ``If God permits it, then God must have intended it,'' so they say.

I need to read Laurus soon.

\end{sffamily}
\end{footnotesize}

